\chapter{Introducción}

El presente informe expone el diseño analítico, la verificación mediante simulación y la comprobación experimental de un
circuito de acondicionamiento de señales analógicas de campo dentro de un lazo industrial de corriente y dos circuitos de
filtros activos de orden superior empleando diversas aproximaciones matemáticas y topologías circuitales, con un tercer
circuito de filtro que no fue implementado.

\section{Objetivos Generales}

\begin{itemize}
    \item Diseñar y caracterizar estática y dinámicamente un circuito de acondicionamiento de señal para un bucle de $4
        \unit{\milli\ampere}$ a $20 \unit{\milli\ampere}$ alimentado con fuente de $+24 \unit{\volt}$.
    \item Sintetizar tres filtros activos con parámetros asignados: un pasa-bajos de 4to orden (Butterworth, Sallen-Key,
        $f_p = 19 \unit{\kilo\hertz}$, $A_v = 9 \unit{\volt\per\volt}$), un pasa-altos de 4to orden (Chebyshev, MFB,
        $f_p = 11 \unit{\kilo\hertz}$, $A_v = 6 \unit{\volt\per\volt}$) y un pasa-banda de 2do orden (Butterworth,
        Sallen-Key, $f_o = 25 \unit{\kilo\hertz}$, $BW = 2 \unit{\kilo\hertz}$, $A_v = 4 \unit{\volt\per\volt}$).
    \item Implementar físicamente en circuito impreso dos de las topologías diseñadas.
    \item Relevar experimentalmente en el laboratorio la ganancia en continua, la respuesta temporal y los perfiles de
        densidad espectral (FFT) en osciloscopio digital, contrastando los datos empíricos contra las predicciones
        teóricas y las simulaciones en SPICE.
\end{itemize}
